\section{CPU memory systems}
\label{sec:cpu-memory-systems}

\subsection{Caches, main memory and physical residency}

A CPU load or store uses a virtual address that is translated to a
physical page.  Data is transferred through a hierarchy of registers,
CPU caches and main memory.  Caches retain recently used cache lines,
not whole model objects as an indivisible unit.  A cache hit can serve
an access without fetching the line from main memory; a miss requires
data from a lower level.  Cache levels differ in capacity, latency and
sharing, and the exact organization is processor dependent.

The \textbf{working set} is the data used during the interval under
study.  Cache capacity can allow a tile of a large matrix to be reused
even when the complete model is much larger than the last-level cache.
Conversely, a tensor's presence in RAM does not imply a cache hit.
Stores also generate traffic: modified cache lines eventually need
writeback, and some store paths fetch a line before updating it.  The
bytes logically consumed by a kernel are consequently different from
the bytes observed at a memory controller.

File size describes persistent storage.  Virtual address space describes
addressable ranges.  Allocated capacity describes storage assigned to
objects or pools.  \textbf{Physical residency} describes which pages are
currently backed in RAM.  Operating-system resident-set measurements
must be interpreted with care when pages are shared or file backed;
summing process footprints can count shared pages more than once.
These quantities answer different questions and should not be used as
interchangeable measures of model memory use.

\subsection{HBM, DDR and memory modes}

\textbf{DDR} is conventional system DRAM accessed through memory
channels.  \textbf{HBM} uses stacked DRAM and a wide interface to provide
high aggregate bandwidth.  In CPU systems combining them, HBM generally
offers a smaller high-bandwidth resource and DDR a larger capacity
resource; usable sizes and attainable rates depend on the platform.
HBM's bandwidth advantage does not establish a lower latency for every
access pattern~\cite{intelmaxoverview,ibeid2025aurora}.

\textbf{Bandwidth} is transferred bytes per unit time,
\textbf{latency} is the time to complete an access, and
\textbf{capacity} is the amount held at once.  Concurrent independent
misses can hide part of the latency and sustain bandwidth.  A dependent
load chain can remain latency limited even when measured bandwidth is
low.  A working set can fit in a tier yet fail to use its bandwidth, or
require more capacity than that tier provides despite modest traffic.

In a separately addressable, often called \textbf{Flat}, configuration,
software can allocate physical pages from exposed memory domains.
Placement can then determine which tier backs weight tensors or KV
storage.  In a hardware-managed \textbf{Cache} configuration, HBM acts
as a cache for DDR backing memory; software does not assign every tensor
a separate HBM address range in the same way~\cite{intelmaxoverview}.
\figref{fig:cpu-memory-organization} shows the distinction.  Neither
mode is a universal default; the actual firmware mode and exposed
memory topology belong to the platform description.

\begin{figure}[htbp]
\centering
% Flat and Cache mode only; allocation choices and access paths are distinct.
\begingroup
\newcommand{\modecore}[1]{%
 \begin{scope}[shift={(#1,0.65)}]
  \foreach \p in {-0.4,0,0.4} {
   \draw[draw=black!45] (\p,-0.55)--(\p,-0.7);
   \draw[draw=black!45] (\p,0.55)--(\p,0.7);
   \draw[draw=black!45] (-0.65,\p)--(-0.8,\p);
   \draw[draw=black!45] (0.65,\p)--(0.8,\p);
  }
  \draw[hardware,rounded corners=2pt] (-0.65,-0.55) rectangle (0.65,0.55);
  \foreach \x in {-0.4,0.05} \foreach \y in {-0.32,0.1}
   \draw[draw=tabline,fill=white] (\x,\y) rectangle ++(0.35,0.24);
 \end{scope}
 \node[note,anchor=west] at (#1+1.0,0.65) {Cores};
 \begin{scope}[shift={(#1,-0.95)}]
  \draw[hardware,rounded corners=2pt] (-0.9,-0.4) rectangle (0.9,0.4);
  \foreach \x in {-0.68,-0.23,0.22} \foreach \y in {-0.25,0.08}
   \draw[draw=tabline,fill=white] (\x,\y) rectangle ++(0.37,0.2);
 \end{scope}
 \node[note,anchor=west] at (#1+1.08,-0.95) {CPU caches};
 \draw[access] (#1,-0.09)--(#1,-0.53);
}
\newcommand{\modestack}[2]{%
 \begin{scope}[shift={(#1,#2)}]
  \foreach \dy in {0,0.22,0.44}
   \draw[draw=hbmcol,fill=hbmcol!12,line width=0.65pt]
    (-0.8,\dy)--(0.45,\dy)--(0.9,\dy+0.35)
    --(-0.35,\dy+0.35)--cycle;
  \foreach \dx in {-0.25,0.15,0.55}
   \draw[draw=hbmcol!60] (\dx,0.46)--++(0.24,0.18);
 \end{scope}
}
\newcommand{\modedimm}[2]{%
 \begin{scope}[shift={(#1,#2)}]
  \draw[draw=ddrcol,fill=ddrcol!8,rounded corners=1pt,line width=0.65pt]
   (-1.1,-0.4) rectangle (1.1,0.4);
  \foreach \dx in {-0.85,-0.2,0.45}
   \draw[draw=ddrcol,fill=ddrcol!16] (\dx,-0.16) rectangle ++(0.43,0.42);
  \foreach \dx in {-0.85,-0.6,-0.35,-0.1,0.15,0.4,0.65,0.9}
   \draw[draw=ddrcol,line width=1pt] (\dx,-0.4)--(\dx,-0.54);
 \end{scope}
}
\newcommand{\modeobjects}[2]{%
 \begin{scope}[shift={(#1,#2)}]
  \draw[draw=tabline,rounded corners=3pt,line width=0.65pt]
   (-1.48,-0.73) rectangle (1.48,0.55);
  \draw[draw=figcol,fill=figcol!8] (-1.05,-0.12) rectangle (-0.35,0.42);
  \foreach \dx in {-0.82,-0.58}
   \draw[draw=figcol!50] (\dx,-0.12)--(\dx,0.42);
  \draw[draw=figcol!50] (-1.05,0.15)--(-0.35,0.15);
  \foreach \dx in {0.38,0.62,0.86} \foreach \dy in {-0.1,0.18}
   \draw[draw=ddrcol,fill=ddrcol!15] (\dx,\dy) rectangle ++(0.19,0.22);
  \node[note] at (-0.7,-0.45) {Weights};
  \node[note] at (0.72,-0.45) {KV};
 \end{scope}
}
\resizebox{0.92\textwidth}{!}{%
\begin{tikzpicture}[
 font=\small\sffamily,text=black!80,
 hardware/.style={draw=tabline,fill=neutralcol,line width=0.65pt},
 access/.style={-{Stealth[length=4pt,width=3pt]},draw=black!65,line width=0.85pt},
 placement/.style={-{Stealth[length=4pt,width=3pt]},draw=black!55,dashed,line width=0.75pt},
 note/.style={font=\footnotesize\sffamily,align=center},
 heading/.style={font=\small\sffamily\bfseries,align=center}
]
\node[heading] at (2,2.05) {(a) Flat};
\node[note,text=black!60] at (2,1.56) {Software-visible tiers};
\node[heading] at (9.8,2.05) {(b) Cache};
\node[note,text=black!60] at (9.8,1.56) {Hardware-managed HBM};
\modecore{2}
\modecore{9.8}
\draw[draw=tabline,dashed] (5.9,1.85)--(5.9,-6.5);

% Flat: two separate access paths, with no default tensor-to-tier assignment.
\modestack{0}{-2.85}
\modedimm{4}{-2.5}
\node[note] at (0,-3.4) {HBM};
\node[note] at (4,-3.4) {DDR};
\draw[access] (2,-1.37)--(2,-1.73)--(0,-1.73)--(0,-2.02);
\draw[access] (2,-1.73)--(4,-1.73)--(4,-2.06);
\fill[black!65] (2,-1.73) circle (1.1pt);
\modeobjects{2}{-4.9}
\node[note,text=black!60] at (2,-4.05) {Backing-page placement};
% Route alternatives outside the memory symbols and their labels.
\draw[placement] (0.5,-4.75)--(-1.15,-4.75)--(-1.15,-2.6)--(-0.8,-2.6);
\draw[placement] (3.5,-4.75)--(5.35,-4.75)--(5.35,-2.5)--(5.12,-2.5);

% Cache: CPU caches are distinct from the memory-side HBM cache.
\modestack{9.8}{-2.85}
\modedimm{9.8}{-4.25}
\node[note,anchor=west] at (10.92,-2.45) {HBM cache};
\node[note,anchor=west] at (11.1,-4.25) {DDR};
\draw[access] (9.8,-1.37)--(9.8,-2.02);
\draw[access] (9.8,-2.87)--(9.8,-3.82);
\node[note,anchor=west,text=black!60] at (10.23,-3.4) {HBM miss};
\modeobjects{9.8}{-5.8}
\draw[placement] (9.8,-5.23)--(9.8,-4.81);

\end{tikzpicture}%
}
\endgroup

\caption{Flat and Cache memory organizations.  In Flat mode, backing
pages for each weight or KV allocation can reside in either exposed
tier.  Dashed arrows show possible backing storage, without choosing a
placement policy.  In Cache mode, allocations have DDR backing storage;
hardware manages the HBM cache, and an HBM miss accesses DDR.  CPU
caches form a separate hierarchy in both modes.  Solid arrows show
access paths, not all individual transfers.  Symbols and sizes are
schematic; stacked dies represent HBM and a module represents DDR.}
\label{fig:cpu-memory-organization}
\end{figure}
\FloatBarrier

\subsection{Sockets, NUMA and allocation locality}

A shared-memory node can contain multiple processor sockets.  With
\textbf{non-uniform memory access} (NUMA), memory domains have different
access costs depending on which cores issue requests.  A remote access
can traverse an inter-socket link and compete for resources beyond the
local memory channels.  A socket can also expose several locality
domains.  NUMA is therefore a relationship between initiator cores and
memory targets, not simply a synonym for socket count~\cite{linuxnumaperf}.
\figref{fig:numa-access-paths} contrasts local and remote access from
the same executing core.

\begin{figure}[htbp]
\centering
% One initiator accesses local or remote pages; domains are not fixed to sockets.
\begingroup
\resizebox{0.92\textwidth}{!}{%
\begin{tikzpicture}[
 font=\small\sffamily,text=black!80,
 hardware/.style={draw=tabline,fill=neutralcol,line width=0.65pt},
 local/.style={-{Stealth[length=5pt,width=4pt]},draw=figcol,line width=1.2pt},
 remote/.style={-{Stealth[length=5pt,width=4pt]},draw=ddrcol,dashed,line width=1.2pt},
 note/.style={font=\footnotesize\sffamily,align=center},
 heading/.style={font=\small\sffamily\bfseries}
]
\foreach \xmin/\xmax/\mid/\name in {0/5/2.5/A,7.3/12.3/9.8/B} {
 \draw[draw=tabline,rounded corners=5pt,line width=0.8pt]
  (\xmin,0) rectangle (\xmax,-3.6);
 \node[heading] at (\mid,0.45) {NUMA domain \name};
}

% Core symbols; the first core is the highlighted access initiator.
\foreach \cx in {1.2,8.5} {
 \begin{scope}[shift={(\cx,-1.4)}]
  \foreach \p in {-0.38,0,0.38} {
   \draw[draw=black!45] (\p,-0.65)--(\p,-0.8);
   \draw[draw=black!45] (\p,0.65)--(\p,0.8);
   \draw[draw=black!45] (-0.65,\p)--(-0.8,\p);
   \draw[draw=black!45] (0.65,\p)--(0.8,\p);
  }
  \draw[hardware,rounded corners=2pt] (-0.65,-0.65) rectangle (0.65,0.65);
  \foreach \dx in {-0.4,0.05} \foreach \dy in {-0.4,0.05}
   \draw[draw=tabline,fill=white] (\dx,\dy) rectangle ++(0.35,0.35);
 \end{scope}
}
\draw[draw=figcol,fill=figcol!25] (1.02,-1.64)--(1.02,-1.16)--(1.43,-1.4)--cycle;
\node[note] at (1.2,-0.35) {Executing core};
\node[note] at (8.5,-0.35) {Other cores};

% Generic physical pages: the comparison is locality, not HBM versus DDR.
\foreach \cx/\name in {3.85/A,11.15/B} {
 \begin{scope}[shift={(\cx,-1.4)}]
  \draw[hardware,fill=white] (-0.42,-0.38) rectangle (0.78,0.78);
  \draw[hardware,fill=white] (-0.6,-0.56) rectangle (0.6,0.6);
  \foreach \dx in {-0.43,0.05} \foreach \dy in {-0.37,0.08}
   \draw[hardware] (\dx,\dy) rectangle ++(0.38,0.34);
 \end{scope}
 \node[note] at (\cx,-2.58) {Pages in \name};
}

% Local requests stay in A. Remote requests traverse a domain interconnect.
\draw[local] (2.03,-1.4)--(3.22,-1.4);
\node[note] at (2.63,-0.98) {Local};
\draw[hardware,fill=white,rounded corners=2pt]
 (5.36,-3.39) rectangle (6.94,-2.81);
\foreach \dy in {-3.24,-2.96}
 \draw[draw=tabline] (5.5,\dy)--(6.8,\dy);
\draw[remote] (1.2,-2.21)--(0.35,-2.21)--(0.35,-3.1)
 --(10.1,-3.1)--(10.1,-1.4)--(10.52,-1.4);
\node[note,fill=white,inner sep=2pt] at (3.45,-3.1) {Remote};
\node[note] at (6.15,-3.82) {Interconnect};
\end{tikzpicture}%
}
\endgroup

\caption{Local and remote NUMA access from a core in domain A.  Local
access targets pages in A; remote access reaches pages in B through an
interconnect.  The additional path and resource contention can change
latency and attainable bandwidth, without a universal numerical penalty.
Domains can lie within one socket or across sockets.  Arrows indicate
access requests; CPU caches and memory controllers are omitted to focus
on locality.  Page symbols do not identify a particular memory tier.}
\label{fig:numa-access-paths}
\end{figure}
\FloatBarrier

\textbf{Thread affinity} determines where threads execute; memory policy
influences which domains provide pages.  They are separate controls.
For ordinary anonymous allocations under a local default policy,
physical pages are commonly established on first touch rather than when
a virtual range is merely reserved.  Parallel initialization can thus
produce different locality from serial initialization.  Binding or
interleaving policies change allocation behavior, and file-backed mappings
have their own policy rules~\cite{linuxnumapolicy2026}.  Touching a tensor
on one core does not guarantee that later accesses by all other cores
are local.  Observed page placement and the loading mechanism must
support such a claim.

\textbf{Placement} describes where backing storage resides.
\textbf{Access traffic} describes bytes transferred while using it.
\textbf{Migration} changes existing data's residency, for example by
moving pages between domains.  Allocating a destination and copying
into it creates additional traffic and may temporarily require both
copies.  Cache-line fills do not by themselves constitute migration of
the tensor's backing pages.  These distinctions apply to both weight
and KV storage, independent of which policies a later design examines.

A cluster connects independently booted nodes and allocates them
through a scheduler.  Shared-memory access within a node differs from
explicit communication between nodes.  The actual topology, scheduler
workflow and operating conditions are developed in the platform and
methodology material.
\FloatBarrier
